\section{总结与延伸}
Lecture 5 把 GPU 性能拆成三个层级：硬件结构、算子组织、算法重排。硬件层给出 SM/warp/shared memory/HBM/tensor cores；算子层使用 low precision、fusion、coalescing、tiling；算法层用 FlashAttention 这类 IO-aware 设计避免不必要 HBM 写回。
\begin{importantbox}{最终 takeaway}
大模型训练不是“把 PyTorch 放到 GPU 上”这么简单。真正的 GPU-aware thinking 是：每个 byte 从 HBM 进来后要被复用多少次？中间结果是否必须写回？thread/warp/block 是否对齐硬件？算法能否重排成更少 IO 的形式？
\end{importantbox}
\subsection{拓展阅读}
\begin{itemize}
    \item NVIDIA CUDA Programming Guide and H100/B200 architecture material.
    \item JAX Scaling Book roofline and inference chapters.
    \item Dao et al. FlashAttention papers.
    \item NVIDIA Transformer Engine low precision documentation.
\end{itemize}
\end{document}
